% =============================================================================
% HILTI SMART CHARGING — Verhaltensorientierte Pricing-Strategie
% EBF Session: EBF-S-2026-02-17-ORG-001
% Datum: 17. Februar 2026
% Modus: STANDARD
% Format: LaTeX 10-Pager (8D: D1=0.7, D4=0.6, D5=G2, D8=0.7)
% =============================================================================

\documentclass[11pt,a4paper]{article}

% --- Packages ---
\usepackage[utf8]{inputenc}
\usepackage[T1]{fontenc}
\usepackage[ngerman]{babel}
\usepackage{geometry}
\usepackage{graphicx}
\usepackage{booktabs}
\usepackage{longtable}
\usepackage{amsmath,amssymb}
\usepackage{xcolor}
\usepackage{hyperref}
\usepackage{fancyhdr}
\usepackage{titlesec}
\usepackage{enumitem}
\usepackage{tcolorbox}
\usepackage{multirow}
\usepackage{float}

% --- FehrAdvice Colors ---
\definecolor{fadarkblue}{RGB}{2,64,121}
\definecolor{falightblue}{RGB}{84,158,222}
\definecolor{fadarkgray}{RGB}{37,33,42}
\definecolor{falightgray}{RGB}{243,245,247}
\definecolor{faorange}{RGB}{222,157,62}
\definecolor{famint}{RGB}{126,189,172}
\definecolor{fared}{RGB}{200,60,60}

% --- Page Setup ---
\geometry{top=2.5cm, bottom=2.5cm, left=2.5cm, right=2.5cm}
\setlength{\parindent}{0pt}
\setlength{\parskip}{6pt}

% --- Header/Footer ---
\pagestyle{fancy}
\fancyhf{}
\fancyhead[L]{\small\color{fadarkblue}FehrAdvice \& Partners AG}
\fancyhead[R]{\small\color{fadarkgray}Hilti Smart Charging Strategy}
\fancyfoot[C]{\small\color{fadarkgray}\thepage}
\fancyfoot[R]{\small\color{fadarkgray}Version 1.0 | 17.02.2026}
\renewcommand{\headrulewidth}{0.5pt}

% --- Section Formatting ---
\titleformat{\section}{\Large\bfseries\color{fadarkblue}}{}{0em}{}
\titleformat{\subsection}{\large\bfseries\color{fadarkblue}}{}{0em}{}
\titleformat{\subsubsection}{\normalsize\bfseries\color{fadarkgray}}{}{0em}{}

% --- Hyperlinks ---
\hypersetup{colorlinks=true, linkcolor=fadarkblue, urlcolor=falightblue, citecolor=fadarkgray}

% --- Custom Boxes ---
\newtcolorbox{keyinsight}{
  colback=falightgray, colframe=fadarkblue,
  title={\textbf{Kernerkenntnis}}, fonttitle=\color{white},
  colbacktitle=fadarkblue, boxrule=0.5pt
}

\newtcolorbox{warningbox}{
  colback=fared!5, colframe=fared,
  title={\textbf{Warnung}}, fonttitle=\color{white},
  colbacktitle=fared, boxrule=0.5pt
}

\newtcolorbox{recommendbox}{
  colback=famint!10, colframe=famint!80!black,
  title={\textbf{Empfehlung}}, fonttitle=\color{white},
  colbacktitle=famint!80!black, boxrule=0.5pt
}

% =============================================================================
\begin{document}

% --- Title Page ---
\begin{titlepage}
\centering
\vspace*{3cm}

{\color{fadarkblue}\rule{\textwidth}{2pt}}
\vspace{1cm}

{\Huge\bfseries\color{fadarkblue} Hilti Smart Charging}\\[0.5cm]
{\LARGE\color{fadarkgray} Verhaltensorientierte Pricing-Strategie\\f\"ur E-Ladestationen am Arbeitsplatz}

\vspace{1.5cm}

{\large\color{fadarkblue}\rule{0.5\textwidth}{0.5pt}}

\vspace{1cm}

{\large\color{fadarkgray}
Evidence-Based Framework (EBF) Analyse\\[0.3cm]
Session: EBF-S-2026-02-17-ORG-001\\[0.3cm]
Modus: STANDARD
}

\vspace{2cm}

{\large
\begin{tabular}{ll}
\textbf{\color{fadarkblue}Auftraggeber:} & Hilti AG, Schaan (LI) \\[0.2cm]
\textbf{\color{fadarkblue}Erstellt von:} & FehrAdvice \& Partners AG \\[0.2cm]
\textbf{\color{fadarkblue}Datum:} & 17. Februar 2026 \\[0.2cm]
\textbf{\color{fadarkblue}Version:} & 1.0 \\
\end{tabular}
}

\vfill
{\color{fadarkblue}\rule{\textwidth}{2pt}}

{\small\color{fadarkgray} FehrAdvice \& Partners AG | Prof. Ernst Fehr | Evidence-Based Framework}
\end{titlepage}

% --- Table of Contents ---
\tableofcontents
\newpage

% =============================================================================
\section{Executive Summary}
% =============================================================================

Hilti betreibt 100 E-Ladestationen f\"ur ca. 400 Mitarbeitende --- ein Verh\"altnis von 4:1. Das Kernproblem ist nicht fehlende Infrastruktur, sondern \textbf{Blockier-Verhalten}: Fahrzeuge belegen Ladestationen stundenlang nach Abschluss des Ladevorgangs.

\begin{keyinsight}
Eine einfache Blockier-Geb\"uhr ist die \textbf{falscheste L\"osung}. Sie w\"urde die bestehende Workplace-Norm \guillemotleft Sei r\"ucksichtsvoll\guillemotright\ durch eine Markt-Norm \guillemotleft Ich zahle daf\"ur\guillemotright\ ersetzen --- \textbf{irreversibel}. Dieses Ph\"anomen (Crowding-Out, $\gamma = -0.68$) ist empirisch umfassend dokumentiert (Gneezy \& Rustichini 2000, Frey \& Jegen 2001).
\end{keyinsight}

Stattdessen empfehlen wir eine \textbf{3-Schichten-Architektur}:

\begin{enumerate}[leftmargin=2cm]
\item[\textbf{Schicht 1:}] \textbf{Awareness} (Woche 1--2) --- Notifications l\"osen $\sim$50\% des Problems
\item[\textbf{Schicht 2:}] \textbf{Social Norms} (Woche 3--6) --- Team-Rankings l\"osen weitere $\sim$30\%
\item[\textbf{Schicht 3:}] \textbf{Strukturelles Pricing} (ab Woche 7) --- Zeitslots oder Deposit f\"ur Rest
\end{enumerate}

\begin{table}[H]
\centering
\rowcolors{2}{falightgray}{white}
\begin{tabular}{lrr}
\toprule
\textbf{\color{fadarkblue}Kennzahl} & \textbf{\color{fadarkblue}Status Quo} & \textbf{\color{fadarkblue}Nach Intervention} \\
\midrule
Effektive Kapazit\"at (Autos/Tag) & 38 & 172 \\
Abdeckung (von 400 MA) & 9\% & 43\% \\
$\varnothing$ Blockier-Zeit nach Ladung & 5.0 h & 0.5 h \\
Investition & --- & CHF 8'000--10'000 \\
Alternative (300 neue Stationen) & --- & CHF 3'000'000+ \\
\bottomrule
\end{tabular}
\caption{Zusammenfassung: Erwartete Wirkung der 3-Schichten-Strategie}
\end{table}


% =============================================================================
\section{Problemdefinition}
% =============================================================================

\subsection{Ausgangslage}

Hilti stellt 100 E-Ladestationen f\"ur Mitarbeitende zur Verf\"ugung. Bei ca. 400 Mitarbeitenden mit Elektrofahrzeug ergibt sich ein Verh\"altnis von 4:1. Eine vollst\"andige Ladung dauert 2--4 Stunden, doch Fahrzeuge verbleiben im Durchschnitt gesch\"atzt 8 Stunden an der Station --- davon 5 Stunden \textbf{nach} Abschluss der Ladung.

\subsection{Das eigentliche Problem: Verhalten, nicht Infrastruktur}

Das Problem ist kein Kapazit\"atsproblem, sondern ein \textbf{Verhaltensproblem}. Die Queue-theoretische Analyse zeigt:

\begin{equation}
c_{\text{eff}} = c \times \frac{t_{\text{charge}}}{t_{\text{charge}} + t_{\text{block}}}
\end{equation}

wobei $c = 100$ (Stationen), $t_{\text{charge}} = 3\text{h}$ (Ladezeit) und $t_{\text{block}}$ (Blockier-Zeit nach Ladung).

\begin{table}[H]
\centering
\rowcolors{2}{falightgray}{white}
\begin{tabular}{lrrr}
\toprule
\textbf{\color{fadarkblue}Szenario} & \textbf{\color{fadarkblue}$t_{\text{block}}$} & \textbf{\color{fadarkblue}$c_{\text{eff}}$/Tag} & \textbf{\color{fadarkblue}2 Durchl\"aufe} \\
\midrule
Status Quo & 5.0 h & 38 & 76 \\
+ Awareness & 2.0 h & 60 & 120 \\
+ Social Norms & 1.0 h & 75 & 150 \\
+ Strukturell & 0.5 h & 86 & \textbf{172} \\
Optimum & 0.25 h & 92 & 184 \\
\bottomrule
\end{tabular}
\caption{Effektive Kapazit\"at als Funktion der Blockier-Zeit}
\end{table}

\begin{keyinsight}
Allein durch Reduktion der Blockier-Zeit von 5h auf 0.5h --- \textbf{ohne eine einzige neue Ladestation} --- steigt die Kapazit\"at um 350\% (von 38 auf 172 Autos/Tag bei 2 Durchl\"aufen).
\end{keyinsight}


% =============================================================================
\section{Kontextanalyse ($\Psi$-Dimensionen)}
% =============================================================================

Das EBF analysiert jeden Kontext \"uber 8 Dimensionen ($\Psi$-Framework, Kapitel 9). F\"ur Hiltis Ladesituation sind 5 Dimensionen besonders relevant:

\subsection{$\Psi_I$ --- Institutioneller Kontext (Regeln)}

\textbf{Status Quo:} Keine formalen Regeln f\"ur Blockier-Verhalten. Der implizite Default ist \guillemotleft unbegrenztes Parken an der Ladestation\guillemotright. Nach der Default-Theorie (PAR-INT-001: $E_{\text{default}} = 0.35$) ist dieser Default verhaltenssteuernd --- er legitimiert Blockieren.

\subsection{$\Psi_S$ --- Sozialer Kontext (Peers)}

\textbf{Workplace als Repeated Game:} Mitarbeitende interagieren t\"aglich \"uber Jahre. Dies erm\"oglicht \textbf{Conditional Cooperation} (MS-SP-008): \guillemotleft Ich mache Platz, wenn andere auch Platz machen.\guillemotright

BCM2-Daten f\"ur die Region (Schweiz/Liechtenstein):
\begin{itemize}[noitemsep]
\item Solidarit\"atsbereitschaft: 75.2\% (CH-SOC-02) --- \textit{hoch}
\item Interpersonales Vertrauen: 61.2\% (CH-SOC-03) --- \textit{\"uberdurchschnittlich}
\item Gemeinwohlorientierung: 7.2/10 (CH-SOC-12) --- \textit{steigend}
\item DE-CH/LI Region: Soziale Normen \textbf{besonders wirksam} (CH-SPR-10)
\end{itemize}

\subsection{$\Psi_C$ --- Kognitiver Zustand}

\textbf{Present Bias} ($\beta = 0.70$, PAR-BEH-003): \guillemotleft Jetzt nicht aufstehen\guillemotright\ wird \"uberbewertet gegen\"uber \guillemotleft sp\"ater Platz machen f\"ur Kolleg:innen\guillemotright.

\textbf{Endowment Effect} (MS-RD-003): Sobald ein:e Mitarbeiter:in an der Station steht, entsteht ein Besitzgef\"uhl (\guillemotleft mein Platz\guillemotright) --- auch wenn der Ladeplatz der Firma geh\"ort.

\subsection{$\Psi_T$ --- Temporaler Kontext}

Arbeitstag-Rhythmus (7--17h) mit Meetings, Deadlines und Fokusphasen. Typischer Ablauf:
\begin{enumerate}[noitemsep]
\item 7:30 --- Ankunft, Auto einstecken
\item 10:30 --- Ladung fertig (3h)
\item 10:30--17:00 --- \textbf{6.5h Blockade} (Meetings, Vergessen, Bequemlichkeit)
\end{enumerate}

\subsection{$\Psi_K$ --- Kultureller Kontext}

Hilti als Premium-Technologieunternehmen in Liechtenstein/Schweiz:
\begin{itemize}[noitemsep]
\item Klimabewusstsein: 65\% der Bev\"olkerung (CH-CLI-04) --- \textit{hoch und steigend}
\item E-Mobility-Anteil: 6.5\% $\rightarrow$ 28\% bis 2030 (CH-INF-04) --- \textit{starkes Wachstum}
\item Nachhaltige Mobilit\"at: 36\% Modal Split (CH-CLI-07) --- \textit{positive Grundhaltung}
\end{itemize}


% =============================================================================
\section{Modellspezifikation}
% =============================================================================

\subsection{Entscheidungsmodell: Bleiben vs. Umparken}

Jede:r Mitarbeiter:in $i$ steht nach Ladeende vor der Entscheidung:

\begin{align}
U_i(\text{bleiben}) &= \underbrace{u_P(\text{convenience})}_{\text{Bequemlichkeit}} - \underbrace{u_F(\text{fee} \times t_{\text{block}})}_{\text{Kosten}} - \underbrace{u_S(\text{guilt} \times \text{visibility})}_{\text{Soziale Kosten}} \\[6pt]
U_i(\text{umparken}) &= -\underbrace{u_P(\text{inconvenience})}_{\text{Aufwand}} + \underbrace{u_S(\text{reputation})}_{\text{Anerkennung}} + \underbrace{u_E(\text{green\_identity})}_{\text{Klima-ID}}
\end{align}

Verhalten: Umparken wenn $U_i(\text{umparken}) > U_i(\text{bleiben})$.

\subsection{Segmentierung: 3 Blocker-Typen}

Nicht alle Blockierer:innen sind gleich. Die Unterscheidung ist \textbf{entscheidend} f\"ur das Interventions-Design:

\begin{table}[H]
\centering
\rowcolors{2}{falightgray}{white}
\begin{tabular}{p{2.5cm}cp{3cm}p{3.5cm}c}
\toprule
\textbf{\color{fadarkblue}Typ} & \textbf{\color{fadarkblue}Anteil} & \textbf{\color{fadarkblue}Ursache} & \textbf{\color{fadarkblue}L\"osung} & \textbf{\color{fadarkblue}$h$-Stock} \\
\midrule
Vergessliche & $\sim$50\% & Awareness $A(\cdot) \approx 0$. Wissen nicht, dass Auto geladen ist & Push-Notification & 0 \\
Bequeme & $\sim$35\% & Present Bias $\beta = 0.70$. Wissen es, aber \guillemotleft gleich...\guillemotright & Social Norms + Anstoss & mittel \\
Strategische & $\sim$15\% & Endowment + keine Konsequenzen. \guillemotleft Mein Recht.\guillemotright & Strukturelle \"Anderung & hoch \\
\bottomrule
\end{tabular}
\caption{Drei Blocker-Typen mit unterschiedlichen Interventions-Ans\"atzen}
\end{table}

\subsection{Habit Stock Dynamik}

Blockier-Verhalten akkumuliert sich als Gewohnheit (MS-TP-009):

\begin{equation}
q_i(t+1) = (1 - \delta_h) \cdot q_i(t) + \text{block}_i(t)
\end{equation}

wobei $q_i(t)$ der Habit-Stock, $\delta_h \approx 0.05$/Woche die Depreziation und $\text{block}_i(t) \in \{0,1\}$.

\subsection{Kritische Interaktion: Crowding-Out}

\begin{warningbox}
\textbf{Das Daycare-Paradox} (Gneezy \& Rustichini 2000): Eine Kindertagesst\"atte f\"uhrte eine Versp\"atungs-Strafe von \$3 ein. Ergebnis: Versp\"atungen \textbf{verdoppelten} sich. Die Strafe ersetzte die moralische Norm durch eine Markt-Norm. Als die Strafe entfernt wurde, blieben die Versp\"atungen hoch --- die Norm war \textbf{irreversibel zerst\"ort}.

\vspace{6pt}
\textbf{EBF-Parameter:} $\gamma(\text{Social}, \text{Financial}) = -0.68$ \quad (PAR-COMP-002, Tier 1)

\vspace{6pt}
\textbf{Konsequenz f\"ur Hilti:} Eine pauschale Blockier-Geb\"uhr transformiert \guillemotleft Ich bin r\"ucksichtsvoll\guillemotright\ in \guillemotleft Ich bezahle f\"ur mein Recht zu blockieren\guillemotright.
\end{warningbox}


% =============================================================================
\section{Parametrisierung}
% =============================================================================

Alle Parameter stammen aus der EBF Parameter-Registry (Layer 2, YAML), nicht aus LLM-Ged\"achtnis:

\begin{table}[H]
\centering
\rowcolors{2}{falightgray}{white}
\begin{tabular}{llccc}
\toprule
\textbf{\color{fadarkblue}ID} & \textbf{\color{fadarkblue}Parameter} & \textbf{\color{fadarkblue}Wert} & \textbf{\color{fadarkblue}CI} & \textbf{\color{fadarkblue}Tier} \\
\midrule
PAR-BEH-001 & $\lambda$ (Loss Aversion) & 2.25 & [1.5, 3.0] & 1 \\
PAR-BEH-003 & $\beta$ (Present Bias) & 0.70 & [0.60, 0.85] & 1 \\
PAR-BEH-002 & $\varphi$ (Crowding-Out) & 0.65 & [0.50, 0.80] & 1 \\
\midrule
PAR-COMP-002 & $\gamma(S,F)$ Social $\times$ Financial & $-$0.68 & [$-$0.78, $-$0.58] & 1 \\
PAR-COMP-001 & $\gamma(\text{ID},S)$ Identity $\times$ Social & $+$0.35 & [0.25, 0.45] & 1 \\
PAR-COMP-004 & $\gamma(S,\text{WG})$ Social $\times$ Warm Glow & $+$0.28 & [0.18, 0.38] & 1 \\
\midrule
PAR-INT-001 & $E_{\text{default}}$ (Default-Effekt) & 0.35 & [0.25, 0.45] & 1 \\
PAR-INT-002 & $E_{\text{social}}$ (Soziale Norm) & 0.020 & [0.015, 0.030] & 1 \\
\bottomrule
\end{tabular}
\caption{EBF-Parameter aus der Parameter-Registry (alle Tier 1 --- Literatur-validiert)}
\end{table}

\subsection{Parameter Context Transformation (PCT)}

Der Labor-Wert $\lambda = 2.25$ wird f\"ur den Hilti-Workplace-Kontext transformiert:

\begin{equation}
\lambda_{\text{Hilti}} = \lambda_{\text{base}} \times M(\Delta\Psi_I) \times M(\Delta\Psi_S) \times M(\Delta\Psi_T)
\end{equation}

\begin{align}
\lambda_{\text{Hilti}} &= 2.25 \times \underbrace{1.3}_{\substack{\text{Gratis-Strom}\\\text{als Referenz}}} \times \underbrace{0.85}_{\substack{\text{Peer}\\\text{Monitoring}}} \times \underbrace{1.1}_{\substack{\text{Arbeitsalltag}\\\text{Zeitdruck}}} \\[6pt]
&= 2.25 \times 1.3 \times 0.85 \times 1.1 \\[6pt]
&\approx \mathbf{2.73}
\end{align}

\textbf{Interpretation:} Loss Aversion ist im Workplace \textbf{h\"oher} als im Labor ($\lambda = 2.73$ vs. 2.25), weil der Verlust von Gratis-Strom als gr\"osserer Loss empfunden wird. Das Peer-Monitoring reduziert leicht (soziale Kosten des Blockierens).


% =============================================================================
\section{Empfohlene Strategie: 3-Schichten Smart Charging}
% =============================================================================

\begin{recommendbox}
\textbf{Kernprinzip:} Nicht \guillemotleft eine L\"osung f\"ur alle\guillemotright, sondern \textbf{sequenzielle Schichten}, die je ein Segment adressieren und dabei die soziale Norm \textbf{sch\"utzen}.
\end{recommendbox}

% --- Schicht 1 ---
\subsection{Schicht 1: Awareness System (Woche 1--2)}

\begin{table}[H]
\centering
\begin{tabular}{ll}
\toprule
\textbf{\color{fadarkblue}Feld} & \textbf{\color{fadarkblue}Wert} \\
\midrule
10C-Target & AWARE (AU) $\rightarrow A(\cdot)\uparrow$ \\
Zielgruppe & Typ 1: Vergessliche ($\sim$50\%) \\
Kosten & $\sim$CHF 2'000 (App-Anpassung) \\
Crowding-Out Risiko & \textcolor{famint}{\textbf{NULL}} \\
Erwarteter Effekt & $-$50\% Blockier-Zeit \\
\bottomrule
\end{tabular}
\end{table}

\textbf{Massnahmen:}
\begin{enumerate}[noitemsep]
\item \textbf{Push-Notification} bei Ladeende: \guillemotleft Dein Auto ist geladen $\checkmark$. Station frei machen?\guillemotright
\item \textbf{+30 Minuten:} \guillemotleft Dein Auto blockiert seit 30 Min. Kolleg:in wartet.\guillemotright
\item \textbf{+60 Minuten:} \guillemotleft Du blockierst seit 1h. Bitte umparken.\guillemotright
\item \textbf{Live-Dashboard} im Eingangsbereich: Freie / Besetzte / Blockierte Stationen
\end{enumerate}

\textbf{Warum zuerst?} Ohne Awareness ($A(\cdot) > 0$) ist \textbf{jede} weitere Intervention wirkungslos. Viele blockieren nicht absichtlich --- sie vergessen es. Die Notification ist das gr\"osste Low-Hanging Fruit.

% --- Schicht 2 ---
\subsection{Schicht 2: Social Norm System (Woche 3--6)}

\begin{table}[H]
\centering
\begin{tabular}{ll}
\toprule
\textbf{\color{fadarkblue}Feld} & \textbf{\color{fadarkblue}Wert} \\
\midrule
10C-Target & WHAT(S) $\rightarrow u_S\uparrow$ und WHO $\rightarrow$ Peer Effects \\
Zielgruppe & Typ 2: Bequeme ($\sim$35\%) \\
Kosten & $\sim$CHF 1'000 (Dashboard-Erweiterung) \\
Crowding-Out Risiko & \textcolor{famint}{\textbf{NULL}} \\
Erwarteter Effekt & Weitere $-$30\% Blockier-Zeit \\
\bottomrule
\end{tabular}
\end{table}

\textbf{Massnahmen:}
\begin{enumerate}[noitemsep]
\item \textbf{Descriptive Norm:} \guillemotleft 87\% deiner Kolleg:innen machen innerhalb von 30 Min. Platz.\guillemotright\ (PAR-INT-002: $E_{\text{social}} = 0.020$)
\item \textbf{Abteilungs-Ranking} (monatlich): \guillemotleft Engineering: $\varnothing$ 22 Min. $|$ Sales: $\varnothing$ 45 Min.\guillemotright
\item \textbf{Green Charger Badge:} $\varnothing$ Blockierzeit $<$ 20 Min. $\rightarrow$ sichtbare Anerkennung
\item \textbf{Management geht voran:} CEO/GL laden und bewegen sichtbar ($\gamma(\text{ID},S) = +0.35$)
\end{enumerate}

\textbf{Theoretische Fundierung:}
\begin{itemize}[noitemsep]
\item MS-SP-008 (Conditional Cooperation): Workplace = repeated game $\rightarrow$ Kooperation stabil wenn sichtbar
\item MS-IN-005 (Social Norms): Descriptive + Injunctive Norm verst\"arken sich
\item PAR-COMP-001: $\gamma(\text{Identity}, \text{Social}) = +0.35$ --- \textbf{positive Synergie!}
\end{itemize}

% --- Schicht 3 ---
\subsection{Schicht 3: Strukturelles Pricing (ab Woche 7)}

\begin{table}[H]
\centering
\begin{tabular}{ll}
\toprule
\textbf{\color{fadarkblue}Feld} & \textbf{\color{fadarkblue}Wert} \\
\midrule
10C-Target & WHEN(V) $\rightarrow \Psi_I$ (Defaults) + READY(AV) $\rightarrow$ Pre-Commitment \\
Zielgruppe & Typ 3: Strategische ($\sim$15\%) \\
Kosten & $\sim$CHF 5'000 (Buchungssystem) \\
Crowding-Out Risiko & \textcolor{famint}{\textbf{GERING}} (bei korrektem Design) \\
Erwarteter Effekt & Verbleibende $-$15--20\% Blockier-Zeit \\
\bottomrule
\end{tabular}
\end{table}

\textbf{NUR einf\"uhren wenn Schicht 1+2 nicht ausreichen!}

\subsubsection{Option A: Zeitslot-Buchungssystem (empfohlen)}

\begin{itemize}[noitemsep]
\item Morgens buchen: \guillemotleft Ich lade 8:00--11:00\guillemotright\ (3h Slot)
\item Automatische Erinnerung 30 Min. vor Slot-Ende
\item \textbf{Default = Slot endet} $\rightarrow$ Status Quo Bias arbeitet F\"UR Umparken
\item P\"unktlich umgeparkt $\rightarrow$ Priority Booking am n\"achsten Tag
\item $2\times$ \"uberziehen $\rightarrow$ 48h Buchungs-Sperre
\item \textbf{Kein monet\"arer Preis} $\rightarrow$ \textbf{kein Crowding-Out}
\end{itemize}

\textbf{Wirkmechanismus:}
\begin{itemize}[noitemsep]
\item MS-TP-007 (Pre-Commitment): Zukunfts-Ich bindet Gegenw\"arts-Ich
\item PAR-INT-001 ($E_{\text{default}} = 0.35$): Default = Slot-Ende, nicht unbegrenztes Parken
\item MS-RD-004 (Status Quo Bias): Jetzt arbeitet der Bias F\"UR das gew\"unschte Verhalten
\end{itemize}

\subsubsection{Option B: Deposit-Refund System (Alternative)}

\begin{itemize}[noitemsep]
\item CHF 10 Deposit beim Einstecken (automatisch via Badge/App)
\item $\leq$30 Min. nach Ladung umgeparkt $\rightarrow$ CHF 10 zur\"uck
\item 30--60 Min. $\rightarrow$ CHF 5 zur\"uck
\item $>$60 Min. $\rightarrow$ CHF 0 zur\"uck
\item Framing: \guillemotleft Hol dir dein Geld zur\"uck\guillemotright\ (\textbf{Gain-Frame}, nicht Loss-Frame)
\item Einnahmen $\rightarrow$ Nachhaltigkeitsfonds ($\gamma(S,\text{WG}) = +0.28$)
\end{itemize}

\textbf{Warum Deposit-Refund statt Strafe?}

\begin{table}[H]
\centering
\rowcolors{2}{falightgray}{white}
\begin{tabular}{p{5cm}p{5cm}}
\toprule
\textbf{\color{fared}Strafe (FALSCH)} & \textbf{\color{famint}Deposit-Refund (RICHTIG)} \\
\midrule
Referenzpunkt = \guillemotleft Null Kosten\guillemotright & Referenzpunkt = \guillemotleft CHF 10 bereits bezahlt\guillemotright \\
Fee = LOSS $\rightarrow$ $\lambda \times$ Fee & Refund = GAIN bei richtigem Verhalten \\
\guillemotleft Ich zahle, also darf ich\guillemotright & \guillemotleft Ich hole MEIN Geld zur\"uck\guillemotright \\
Norm zerst\"ort ($\gamma = -0.68$) & Norm bleibt intakt \\
\bottomrule
\end{tabular}
\caption{Strafe vs. Deposit-Refund: Framing-Unterschied}
\end{table}


% =============================================================================
\section{Sequenzielle Einf\"uhrung}
% =============================================================================

\begin{keyinsight}
Die sequenzielle Einf\"uhrung ist \textbf{nicht optional}. Sie ist wissenschaftlich begr\"undet: (1) Misst den kausalen Effekt jeder Schicht. (2) Vermeidet Crowding-Out, weil die Norm sich \textbf{vor} dem Pricing etabliert. (3) Spart Kosten --- vielleicht reichen Schicht 1+2.
\end{keyinsight}

\begin{table}[H]
\centering
\rowcolors{2}{falightgray}{white}
\begin{tabular}{cllcc}
\toprule
\textbf{\color{fadarkblue}Phase} & \textbf{\color{fadarkblue}Zeitraum} & \textbf{\color{fadarkblue}Massnahme} & \textbf{\color{fadarkblue}Messung} & \textbf{\color{fadarkblue}Ziel} \\
\midrule
1 & Woche 1--2 & Awareness (Notifications) & $\varnothing t_{\text{block}}$ & $\leq$ 2.5h \\
--- & Woche 2 & \textit{Auswertung Phase 1} & Effekt messen & --- \\
2 & Woche 3--6 & + Social Norms & $\varnothing t_{\text{block}}$ & $\leq$ 1.0h \\
--- & Woche 6 & \textit{Auswertung Phase 1+2} & Effekt messen & --- \\
3 & Ab Woche 7 & + Zeitslots/Deposit & $\varnothing t_{\text{block}}$ & $\leq$ 0.5h \\
--- & Woche 12 & \textit{Gesamtauswertung} & Alle KPIs & --- \\
\bottomrule
\end{tabular}
\caption{Implementierungs-Roadmap mit Messzeitpunkten}
\end{table}

\textbf{Entscheidungsregel nach Phase 2:}
\begin{itemize}[noitemsep]
\item $\varnothing t_{\text{block}} \leq 1.0$h $\rightarrow$ Schicht 3 \textbf{nicht n\"otig} (Kosten gespart!)
\item $\varnothing t_{\text{block}} > 1.0$h $\rightarrow$ Schicht 3 aktivieren (Zeitslot bevorzugt)
\end{itemize}


% =============================================================================
\section{Sensitivit\"atsanalyse}
% =============================================================================

\subsection{Welcher Parameter treibt das Ergebnis?}

\begin{table}[H]
\centering
\rowcolors{2}{falightgray}{white}
\begin{tabular}{lcc}
\toprule
\textbf{\color{fadarkblue}Parameter} & \textbf{\color{fadarkblue}Einfluss} & \textbf{\color{fadarkblue}Anteil} \\
\midrule
Awareness $A(\cdot)$ & Haupttreiber & 45\% \\
Social Norm $E_{\text{social}}$ & Stark & 30\% \\
Default $E_{\text{default}}$ (Zeitslot) & Mittel & 15\% \\
Pricing $\lambda \times$ Deposit & Gering & 10\% \\
\bottomrule
\end{tabular}
\caption{Einfluss der Interventions-Komponenten auf Blockier-Reduktion}
\end{table}

\textbf{Robustheit:} Auch ohne Schicht 3 (Pricing) werden 75--80\% des Problems gel\"ost. Das Ergebnis ist robust gegen\"uber Parameter-Unsicherheit, weil der st\"arkste Hebel (Awareness) der am besten dokumentierte ist.

\subsection{Was passiert bei falschem Design?}

\begin{table}[H]
\centering
\rowcolors{2}{falightgray}{white}
\begin{tabular}{p{4cm}p{3cm}p{4cm}}
\toprule
\textbf{\color{fared}Fehler} & \textbf{\color{fared}Effekt} & \textbf{\color{fared}Irreversibel?} \\
\midrule
Flat Fee statt Deposit & Crowding-Out: Blockieren steigt & \textbf{JA} --- Norm zerst\"ort \\
Pricing vor Awareness & Vergessliche werden bestraft & Nein, aber unfair \\
Keine Management-Beteiligung & Norm-Erosion von oben & Teilweise \\
Zu kurze Slots (1h) & Frustration, Umgehung & Nein, anpassbar \\
\bottomrule
\end{tabular}
\caption{Fehler-Szenarien und ihre Konsequenzen}
\end{table}


% =============================================================================
\section{Zusammenfassung: TUN vs. NICHT TUN}
% =============================================================================

\begin{table}[H]
\centering
\begin{tabular}{p{6cm}p{6cm}}
\toprule
\textbf{\color{famint}$\checkmark$ TUN} & \textbf{\color{fared}$\times$ NICHT TUN} \\
\midrule
Notifications einf\"uhren (AWARE ist Voraussetzung) & Flat Fee pro Blockier-Stunde (Daycare-Paradox!) \\[6pt]
Social Norms + Team-Rankings (DE-CH: besonders wirksam) & \guillemotleft Strafzettel\guillemotright\ oder \guillemotleft Bussen\guillemotright\ (Markt-Norm ersetzt Moral-Norm) \\[6pt]
Zeitslot-System mit Default-Ende (Pre-Commitment + Default) & Unbegrenztes Parken mit Fee (\guillemotleft Ich bezahle, also darf ich\guillemotright) \\[6pt]
Deposit-Refund falls Pricing (Gain-Frame) & Alles gleichzeitig einf\"uhren (sequenziell = messbar) \\[6pt]
Sequenziell testen (1$\rightarrow$2$\rightarrow$3) & \guillemotleft Eine L\"osung f\"ur alle\guillemotright\ (3 Segmente = 3 L\"osungen) \\[6pt]
Management geht voran (Signal + Norm) & Ausnahmen f\"ur Gesch\"aftsleitung \\
\bottomrule
\end{tabular}
\caption{Handlungsempfehlungen: Was Hilti tun und nicht tun sollte}
\end{table}


% =============================================================================
\section{Theoretische Fundierung}
% =============================================================================

\begin{table}[H]
\centering
\rowcolors{2}{falightgray}{white}
\begin{tabular}{lll}
\toprule
\textbf{\color{fadarkblue}Theorie-ID} & \textbf{\color{fadarkblue}Name} & \textbf{\color{fadarkblue}Relevanz} \\
\midrule
MS-RD-001 & Prospect Theory (Kahneman \& Tversky 1979) & $\lambda$ wirkt via Deposit \\
MS-RD-003 & Endowment Effect (Thaler 1980) & \guillemotleft Mein Parkplatz\guillemotright \\
MS-RD-004 & Status Quo Bias (Samuelson \& Zeckhauser 1988) & Default = bleiben \\
MS-RD-005 & Mental Accounting (Thaler 1985) & Lade- vs. Blockier-Konto \\
MS-TP-001 & Quasi-Hyperbolic Discounting (Laibson 1997) & Present Bias $\beta$ \\
MS-TP-007 & Pre-Commitment (Thaler \& Benartzi 2004) & Zeitslot-Buchung \\
MS-SP-008 & Conditional Cooperation (Fischbacher et al. 2001) & Workplace als Repeated Game \\
MS-IN-005 & Social Norms (Elster 1989) & Descriptive + Injunctive \\
\midrule
CAS-156 & Tragedy of the Commons & Commons-Problem \\
CAS-005 & Social Norms Feedback for Energy & Norm-Messaging Wirkung \\
CAS-025 & Workplace Reciprocity Norms & Workplace-spezifisch \\
\bottomrule
\end{tabular}
\caption{\"Ubersicht der verwendeten Theorien und Cases}
\end{table}

\subsection{Literaturverweise}

\begin{itemize}[noitemsep]
\item Allcott, H. (2011). Social Norms and Energy Conservation. \textit{Journal of Public Economics}, 95(9-10), 1082--1095.
\item Fischbacher, U., G\"achter, S., \& Fehr, E. (2001). Are people conditionally cooperative? \textit{Economics Letters}, 71(3), 397--404.
\item Frey, B., \& Jegen, R. (2001). Motivation crowding theory. \textit{Journal of Economic Surveys}, 15(5), 589--611.
\item Gneezy, U., \& Rustichini, A. (2000). A Fine is a Price. \textit{Journal of Legal Studies}, 29(1), 1--17.
\item Kahneman, D., \& Tversky, A. (1979). Prospect Theory. \textit{Econometrica}, 47(2), 263--291.
\item Kube, S., Mar\'echal, M.A., \& Puppe, C. (2012). The Currency of Reciprocity. \textit{American Economic Review}, 102(4), 1644--1662.
\item Laibson, D. (1997). Golden Eggs and Hyperbolic Discounting. \textit{Quarterly Journal of Economics}, 112(2), 443--478.
\item Samuelson, W., \& Zeckhauser, R. (1988). Status Quo Bias in Decision Making. \textit{Journal of Risk and Uncertainty}, 1(1), 7--59.
\item Thaler, R.H. (1980). Toward a Positive Theory of Consumer Choice. \textit{Journal of Economic Behavior \& Organization}, 1(1), 39--60.
\item Thaler, R.H., \& Benartzi, S. (2004). Save More Tomorrow. \textit{Journal of Political Economy}, 112(S1), S164--S187.
\end{itemize}


% =============================================================================
\section*{Methodische Anmerkung}
% =============================================================================

{\small
Diese Analyse wurde mit dem \textbf{Evidence-Based Framework (EBF)} erstellt --- einem integrativen Framework f\"ur \"okonomisches und soziales Verhalten basierend auf Komplementarit\"at und Kontext. Alle Parameter stammen aus der EBF Parameter-Registry (Tier 1: Literatur-validiert). Die Kontext-Transformation erfolgt via Parameter Context Transformation (PCT). Keine Werte wurden aus LLM-Ged\"achtnis generiert (Three-Layer Architecture, Prinzip P2: \guillemotleft Parameters from Registry, Not from Memory\guillemotright).

\vspace{6pt}
\textbf{Session:} EBF-S-2026-02-17-ORG-001 \quad \textbf{Modus:} STANDARD \quad \textbf{Version:} 1.0
}

\end{document}
